% ============================================================================
\section{Experimental Setup}
\label{sec:experiments}
% ============================================================================

\subsection{Configuration}

We distill four specialists into a Qwen3-1.7B student in non-thinking mode:
the mathematics and instruction-following teachers are task-specific
Qwen3-1.7B RL checkpoints, and the code and science teachers are Qwen3-4B in
non-thinking mode. A micro-batch contains $b=4$ prompts from one task, one
response per prompt, and one backward pass. Each step accumulates $G=16$
micro-batches, or $64$ responses, with $m_{\min}=2$ and $m_{\max}=8$; the
estimator uses the student's top $k=16$ tokens. Every run lasts $500$ steps,
uses a nonrepeating stream of $16{,}000$ prompts per task, and caps a response
at $4{,}096$ tokens. The fixed objective weights are proportional to the
inverse initial task losses, and timing averages use decay $0.9$. Uniform
assigns four micro-batches to each task, so the count limits imply $H\leq2$.

All training methods use AdamW with \missing{learning rate and schedule,
$\beta_1$, $\beta_2$, $\epsilon$, weight decay, gradient clipping, precision,
and sharding}, and each method uses one seed. A run reserves two GPUs: one
$96$-GB device trains the student and one $48$-GB device generates rollouts
and serves all four teachers. We evaluate $128$ fixed held-out prompts per
task every $50$ steps and use $32$ fresh micro-batches per task at early,
middle, and late checkpoints for the paired variance protocol. Methods share
the prompt order, sampling settings, optimizer, response cap, and checkpoint
schedule. Table~\ref{tab:pipeline} gives the complete configuration and exact
asset revisions.

\subsection{Methods}

The top-$k$ estimator is the default unless a row explicitly says otherwise.
\textbf{Uniform} uses the same count for every task.
\textbf{GPAS (per step)} applies
Equation~\ref{eq:batch_gpas_allocation} to preconditioned noise.
\textbf{GPAS (per GPU hour)} is the same method in its cost-aware mode and
minimizes Equation~\ref{eq:batch_cost_objective}.
\textbf{Counts from unpreconditioned noise} sets $D$ to the identity only when
computing counts; it is the SGD rule run under AdamW and isolates whether the
optimizer's preconditioner matters.
\textbf{Counts from loss gap} assigns counts in proportion to
$w_i\bar L_i$, where $\bar L_i$ is the smoothed training-batch teacher loss,
while retaining the fixed objective.

Two one-component ablations replace the top-$k$ estimator by the usual
sampled-token estimator under, respectively, Uniform counts and GPAS (per
step) counts. These two rows and their top-$k$ counterparts form the estimator
by allocation comparison. \textbf{StdMOPD} uses the sampled-token estimator
and averages all valid tokens in a step together; its effective task weights
therefore follow token counts, so its objective differs from
Equation~\ref{eq:fixed_multitask_objective}. Every adaptive rule begins from a
Uniform step, and all methods otherwise share the configuration above.

\subsection{Measurements}

\paragraph{Held-out teacher loss.}
Held-out teacher loss is the fixed-weight teacher loss on prompts that the
training stream never sees. It is the primary training metric and is evaluated
with a common sampling seed; method differences use paired bootstrap intervals
over prompts, which measure evaluation uncertainty rather than variation over
training seeds. Uniform's final held-out loss is the common target for compute
comparisons.

\paragraph{Gap and token diagnostics.}
At every step and for every task, we log three scalars from the top-$k$ scores.
The gap level is the student-probability-weighted mean log ratio within the
top-$k$ set; the within-position spread is the corresponding weighted variance
of that log ratio; and the tail mass is one minus the summed student probability
of the set. The spread is a low-cost diagnostic rather than the gradient
variance itself, because it omits the score vectors' directions and norms; the
paired gradient measurement below tests whether it predicts token-level noise.

\paragraph{Paired gradient variance.}
At each selected checkpoint, we generate fresh micro-batches and first compute
the actual top-$k$ training gradient, whose sampled tail token comes from the
rollout. These gradients supply $\widehat e_i$ for the held-out allocation
comparison. We split the pool in half: one half estimates task noise and
chooses an allocation, while the other scores that allocation using $V(m)$;
swapping the halves and averaging removes the choice of split. Separately,
after the response traces are fixed, the sampled-token and top-$k$ diagnostic
gradients share an independent scoring-token draw at every prefix. Their paired
variance change isolates the fixed-prefix token mechanism but need not equal
the online training-noise change (details in Appendix~\ref{app:heldout}).

\paragraph{Compute.}
GPU hours are full wall time times the two reserved GPUs, including waiting and
synchronization. A trace assigns count-dependent rollout, scoring, and backward
work to $\tau_i$ and the remaining fixed work to $C$, and we report the fit
residual of this additive model. The overhead ratio is the fixed time per step
divided by time spent on micro-batches; it caps how much a time-aware allocation
can save (details in Appendix~\ref{app:system_details}).

\paragraph{Capability.}
We report MATH-500 greedy pass@1, a fixed LiveCodeBench pass@1 slice, IFBench
strict accuracy, and GPQA-Diamond average@4 for the initial student, each
teacher, and each method. For score $s$, the per-task normalized gain is
$(s-s_{\rm init})/(s_{\rm teacher}-s_{\rm init})$; we also report its average,
mean response length, and truncation rate. \missing{Insert benchmark versions
and complete decoding settings.}

% ============================================================================
\section{Results}
\label{sec:results}
% ============================================================================

\subsection{Noise is heterogeneous, drifting, and coordinate dependent}

\begin{figure}[H]
    \centering
    \missingfigure[0.20\textheight]{Insert Figure 2 from the four-teacher
    logs: (a) preconditioned and unpreconditioned noise over training; (b)
    counts chosen by GPAS in its per-step and per-GPU-hour modes; (c) stacked
    response-level and removable token-level noise by task, aligned with gap
    level and teacher origin.}
    \caption{Task-gradient noise, allocation, and its token component.
    \missing{State the cross-task range, drift, coordinate-ranking disagreement,
    and token share for distant versus same-origin teachers.}}
    \label{fig:allocation_dynamics}
\end{figure}

\begin{table}[H]
\centering
\small
\caption{Held-out preconditioned top-$k$ gradient variance relative to Uniform.
Each allocation is selected on one half of fresh micro-batches and evaluated
on the other. \missing{Fill the checkpoint values and the relative error of
$\widehat e_i$ across halves.}}
\label{tab:heldout_variance}
\begin{tabular}{lccc}
\toprule
Allocation & Early & Middle & Late \\
\midrule
Uniform & $1.00$ & $1.00$ & $1.00$ \\
Counts from unpreconditioned noise & & & \\
Counts from loss gap & & & \\
GPAS (per step) & & & \\
GPAS (per GPU hour) & & & \\
\bottomrule
\end{tabular}
\end{table}

\missing{A single-micro-batch noise estimate spans X across tasks and each
task drifts by up to Y over training; the preconditioned and unpreconditioned
rankings disagree in Z of steps. Report count-bound occupancy and the median
and quartiles of H.} The controlled calculation already shows why the
coordinate choice can matter: the unpreconditioned rule reaches
$2.3\times$ Uniform's preconditioned variance, whereas GPAS reaches
$0.68\times$.

\subsection{Teacher--student gap appears as token-level noise}

Figure~\ref{fig:allocation_dynamics}c separates response-level noise from the
fixed-prefix token component targeted by the top-$k$ estimator. \missing{Report
the token-noise fraction for distant and same-origin teachers, its association
with gap level and within-position spread, the tail mass, and the reduction in
the cross-task noise ratio and resulting count change.} This comparison treats
teacher distance as a measured correlate, not as an assumed monotone law.

\subsection{Analytic removal and sample allocation are complementary}

\begin{table}[H]
\centering
\scriptsize
\setlength{\tabcolsep}{2.5pt}
\caption{Four-teacher capability and efficiency. StdMOPD optimizes a different
objective, so its held-out $F$ and compute-to-target entries are not comparable.
\missing{Populate from evaluation outputs and timing logs; report every
task-level difference from Uniform.}}
\label{tab:main_results}
\begin{tabular}{lccccccc}
\toprule
& Math & Code & IF & Science & Norm. avg. & Final $F$ & GPU h to target \\
\midrule
Initial student & & & & & & & \\
Teachers & & & & & & & \\
StdMOPD & & & & & & & \\
Uniform & & & & & & & \\
GPAS (per step) & & & & & & & \\
GPAS (per GPU hour) & & & & & & & \\
Counts from unpreconditioned noise & & & & & & & \\
Counts from loss gap & & & & & & & \\
Sampled-token estimator (Uniform) & & & & & & & \\
Sampled-token estimator (GPAS per step) & & & & & & & \\
\bottomrule
\end{tabular}
\end{table}

\begin{table}[H]
\centering
\small
\caption{Matched-step $2\times2$ ablation of the gradient estimator and count
allocation. \missing{Fill held-out teacher loss and paired differences from
the four runs.}}
\label{tab:estimator_allocation_ablation}
\begin{tabular}{lcc}
\toprule
& Uniform counts & GPAS (per step) counts \\
\midrule
Sampled-token estimator & & \\
Top-$k$ estimator & & \\
\bottomrule
\end{tabular}
\end{table}

\missing{At matched steps, integrating the retained-token contribution improves
held-out loss by X, allocating the residual noise improves it by Y, and the
combination improves it by Z. Once that contribution is integrated, counts from
loss gap provide no additional advantage over the noise rule.}

\subsection{The overhead ratio bounds the gain from time-aware allocation}

\begin{figure}[H]
    \centering
    \missingfigure[0.17\textheight]{Insert Figure 3: (a) held-out weighted
    teacher loss versus optimizer steps for the matched estimator--allocation
    comparison; (b) the same fixed objective versus GPU hours for Uniform and
    GPAS in its per-step and per-GPU-hour modes. Annotate the measured overhead
    ratio and the time to Uniform's final loss.}
    \caption{Matched-step learning curves and compute efficiency under the
    measured system overhead.
    \missing{State the overhead ratio and GPU-hour saving at the common loss
    target.}}
    \label{fig:training_efficiency}
\end{figure}

\missing{At overhead ratio R, GPAS (per GPU hour) reaches Uniform's final loss
in S fewer GPU hours. Relate the observed saving to the upper limit implied by
fixed step time and report the additive time model's error.}
